Verifiable Institutions, as formalised in the preceding papers of this
series, guarantee algorithmic accountability by encoding the Separation
of Powers into type systems and cryptographic logs. However, this
architecture creates a structural rigidity: the \emph{Dead Hand
problem}. A constitution that cannot change is a prison; a constitution
that changes at the whim of a regulator is susceptible to capture.
Both failure modes are fatal to institutional legitimacy.

This paper introduces the \textbf{Constitutional Amendment Protocol},
a mechanism for formalising institutional evolution within a Verifiable
Institution. We distinguish between \emph{ordinary policy updates}
(parameter changes authorised by the Legislative branch alone) and
\emph{constitutional amendments} (structural changes to the
Separation of Powers invariants), requiring \emph{Tripartite
Ratification}---signatures from all three independent branches---for
the latter. We prove that this threshold is capture-resistant: no
pairwise collusion of two principals suffices to weaken the
fundamental constitutional invariant.

We further encode \emph{Sunset Clauses} as linear resources with
finite lifetime: legislative authority expires unless actively renewed
through a ratification process. This forces a periodic
\emph{Constitutional Interregnum}, ensuring that the algorithm's
authority is re-legitimised by each generation of stakeholders rather
than inherited indefinitely from a prior act.

We prove two theorems---Amendment Soundness and Mandatory Renewal---
and analyse two primary attack vectors: the Hostile Takeover (pairwise
collusion against Stone Clauses) and the Zombie Institution (operating
past mandate expiry). Together, these results complete the framework
for Verifiable Institutions: moving from static \emph{code-as-law}
to dynamic \emph{architecture-as-democracy}.
